\subsection{A free filtered limit forces the monoid comparison}

\begin{claim}{D.8}{Proved}{Product-compatible saturation and the free-limit criterion}
\label{D:free-limit-criterion}
Let $A=\mathcal C_\delta$, and choose its rich countable limit $U$.
Replace $A$ by the equivalent concrete category of nonempty finite
closed substructures of $U$, with all closed embeddings as arrows.
Let $\mathcal I$ be the directed poset of these substructures with
inclusion arrows only, and put $G=\operatorname{Aut}(U)$.

For any monoid diagram $M:A\to\mathrm{Mon}$, form
\[
 S=\colim_{B\in\mathcal I}^{\mathrm{Mon}}M(B),
 \qquad T_S(B)=S^{G_B},
 \qquad G_B=\{g\in G:g|_B=\id_B\}.
 \tag{D.8.1}
\]
The natural map $\eta:M\to T_S$ induces a weak equivalence
$B\int_A M\to B\int_A T_S$ and an isomorphism of ordinary monoid
colimits. More precisely, for every finite $k\geq0$,
\[
 B\El_A(M^k)\longrightarrow B\El_A(T_S^k)
 \quad\text{is a weak equivalence}.
 \tag{D.8.2}
\]
If $M$ takes values in free monoids and $S$ is a free monoid, then
its specified comparison $q_M$ is a weak equivalence.

In particular, let $X:A\to\Set$, $K:A\to\mathrm{Grp}$, and let
$\alpha:F(X)\to K$ be a natural monoid homomorphism, where $F$
denotes the free-monoid functor. The monoids
\[
 R(a)=\{w\in F(X(a)):\alpha_a(w)=1\}
 \tag{D.8.3}
\]
are free, and their diagram always passes its monoid comparison.
This permits arbitrary, including noninjective, transitions in $X$
and $K$.
\end{claim}
\begin{proof}
We use the precise structural inputs established in the all-action
proof of R17, A.30--A.34, audited in D.4
\cite{FTOTEX2}: condition I; strong homogeneity of $U$; directedness
of $\mathcal I$; and contractibility of every nonempty full cosieve
in a slice $B/A$. The last assertion comes from the functorial
ordered free amalgam with one fixed slice object, which receives
natural maps from the identity and from that constant object.
The same construction over the empty base contracts $BA$.
These assertions hold for arbitrary actions, without finiteness or
injectivity assumptions.

\emph{The continuous action and the saturation functor.}
Filtered monoid colimits are computed on underlying sets. Define
\[
 g\cdot[B,m]=[gB,M(g|_B)(m)].
\]
Compatibility with inclusions makes this a well-defined action on
$S$ by monoid automorphisms. Each represented element $[B,m]$ is
fixed by the open subgroup $G_B$, so it is a continuous action on
the discrete monoid $S$. For $f:B\to C$, extend $f$ to $g\in G$
and set $T_S(f)(s)=g\cdot s$. Two extensions differ on the right
by an element of $G_B$, so this is independent of the choice.
Its value belongs to $S^{G_C}$, and extensions of composites prove
functoriality. The maps $m\mapsto[B,m]$ give the natural monoid
map $\eta$.

\emph{Eventual equality and finite products.}
For a Set diagram $H$, put $h\sim_B h'$ if some arrow $f:B\to C$
identifies their images. Condition I makes this an equivalence
relation and makes $H^{\mathrm{inj}}=H/{\sim}$ a functor with
injective transitions. Finitely many equality witnesses can be
synchronized by repeatedly completing two arrows with common source
to a commuting square. Hence injectivization commutes with finite
products. For a monoid diagram the same argument shows that $\sim$
is a monoid congruence.

For the present concrete indexing category, equality in the
filtered colimit over $\mathcal I$ is exactly this eventual
equality. One direction is immediate. For the other, if $f:B\to C$
is a witness, extend it to $g\in G$ and apply the isomorphism
$H(g^{-1}|_C)$. The resulting witness is the inclusion
$B\subseteq g^{-1}C$. Thus $H^{\mathrm{inj}}(B)$ is the image of
$H(B)$ in its filtered colimit $S_H$, and its filtered colimit
is still $S_H$.

For completeness, the two homotopy-final comparisons in R17 apply
as follows. For $q:H\to H^{\mathrm{inj}}$ and an object
$(B,[h])$, its comma category over $\El(q)$ receives a functor
from $B/A$ using the chosen representative $h$. For an object
$((C,k),f)$ of that comma category, the comma over this functor
is the full cosieve of $C/A$ on arrows $u:C\to D$ such that
\[
 H(u)k=H(uf)h.
\]
It is nonempty by the definition of $[k]$ and contractible by
the slice-cosieve contraction. Two applications of Theorem A,
using the initial object of $B/A$, prove that
$B\El(q)$ is a weak equivalence.

For an injective diagram $H\hookrightarrow T_{S_H}$, the comma
over an object $(B,s)$ consists precisely of those arrows
$f:B\to C$ for which $T_{S_H}(f)(s)$ lies in the image of $H(C)$.
The representing element is unique. This is a full cosieve in
$B/A$, and it is nonempty: represent $s$ at some finite closed
$D$ and choose a finite closed $C$ containing $B\cup D$.
Theorem A again gives a weak equivalence of element categories.

Apply these comparisons to $H=M^k$. Injectivization commutes with
finite products by the synchronization argument, filtered colimits
commute with finite products, and taking fixed subsets commutes
with finite products. Consequently these comparisons are exactly
the $k$th powers of the monoid maps above, and prove (D.8.2).
For $k=0$ the map is the identity of the constant singleton action.

\emph{Both colimit comparisons.}
The nerve of the one-object category associated to $M(a)$ has
$k$-simplices $M(a)^k$. Its simplicial replacement in $A$ is
therefore a bisimplicial model for $B\int_A M$ whose rows are
the nerves of the element categories in (D.8.2). Rowwise weak
equivalences induce a weak equivalence on the diagonal. This
proves the asserted comparison of homotopy colimits.

Taking connected components in (D.8.2) gives, separately,
\[
 \colim_A \operatorname{Und}(M)\cong\colim_A \operatorname{Und}(T_S),\qquad
 \colim_A \operatorname{Und}(M^2)\cong\colim_A \operatorname{Und}(T_S^2),
 \tag{D.8.4}
\]
where $\operatorname{Und}$ denotes the underlying-set functor. These bijections
respect both projections, multiplication, and the unit. A monoid
colimit is presented by one generator for each element of
$\colim_A \operatorname{Und}(M)$ and the relations $[x][y]=[xy]$ indexed by
$\colim_A \operatorname{Und}(M^2)$, together with the unit relation. This follows
directly by imposing transition and multiplication relations on
the free monoid on all the stage elements. Thus (D.8.4) identifies
the two monoid-colimit presentations. No interchange of ordinary
Set colimits with products is asserted.

\emph{The free-limit case.}
Suppose $S$ is free, with canonical atom set $E$. Every monoid
automorphism permutes $E$, which is consequently a continuous
discrete $G$-set. A word is fixed by $G_B$ if and only if each
of its ordered letters is fixed. It follows naturally that
\[
                  T_S(B)=F(E^{G_B}).
 \tag{D.8.5}
\]
Thus $T_S$ is a basis-map diagram.

For any Set functor $Y$ on the total category $A$, the comparison
for $F(Y)$ is a weak equivalence. Indeed the bouquet of circles
indexed by $Y(a)$ includes naturally in $BF(Y(a))$ and is an
objectwise homotopy equivalence. Its homotopy colimit is modeled
by attaching $B\El_A Y\times[0,1]$ to $BA$ at both ends via
the projection $B\El_A Y\to BA$. Every component of
$B\El_A Y$ is contractible and $BA$ is contractible. In this
homotopy-pushout description, replacing these spaces by their
components and a point gives a weak equivalence to the bouquet
indexed by $\colim_A Y$. That is the specified comparison to
$BF(\colim_A Y)$, since $F$ preserves colimits.
Apply this observation to $Y(B)=E^{G_B}$. The two comparisons
already proved for $\eta$ form a natural square with $q_M$ and
$q_{T_S}$, proving the free-limit criterion.

\emph{First-return kernels.}
A nonempty word in (D.8.3) is called a first-return word if none
of its proper nonempty prefixes maps to $1$. Cut any return word
at every prefix mapping to $1$. The consecutive pieces are
first-return words, and the cut positions uniquely determine
their factorization. Hence $R(a)$ is free on these words.
An image of a first-return word may split into several first
returns; the resulting maps are nevertheless strict monoid
homomorphisms.

Write $X_\infty=\colim_{\mathcal I}X$ and
$K_\infty=\colim_{\mathcal I}K$. Since filtered colimits commute
with the finite equation $\alpha(w)=1$, and $F$ preserves them,
\[
 \colim_{\mathcal I}R
   =\ker\bigl(F(X_\infty)\longrightarrow K_\infty\bigr).
\]
Explicitly, a finite word and its equality with $1$ are both
represented at a common finite stage; equality of two represented
words is witnessed at another such stage. This proves the equality
on both elements and identifications. Its right-hand monoid is
again free by the first-return argument. The free-limit criterion
therefore applies to $R$ and proves the final assertion.
\end{proof}

The criterion does not assert that every filtered colimit of free
monoids is free. In particular it does not settle the remaining
universal comparison question for $\mathcal C_\delta$.
